\documentclass[11pt]{amsart}
%==============================================================
%  The Bourbaki Engine: a governed discovery architecture for
%  long-horizon autonomous mathematics.
%
%  v2 (discovery reorientation). Integrates the MIRADOR
%  representational substrate and the ProofContext retrieval
%  substrate as bounded modules, under a two-plane
%  (Discovery / Trust) architecture. Companion to the plank
%  pilot (bang-plank-affine-triangle repository).
%==============================================================
\usepackage{amsmath,amssymb,amsthm,mathtools}
\usepackage{booktabs}
\usepackage{array}
\usepackage{longtable}
\usepackage[margin=1in]{geometry}
\usepackage[colorlinks=true,linkcolor=blue,citecolor=blue,urlcolor=blue]{hyperref}
\usepackage{manfnt} % \dbend --- the Bourbaki/Knuth dangerous-bend glyph
\emergencystretch=1.5em

\theoremstyle{definition}
\newtheorem{principle}{Principle}
\newtheorem{definition}{Definition}[section]
\newtheorem{contract}{Contract}[section]
\theoremstyle{remark}
\newtheorem{remark}[definition]{Remark}

\newcommand{\engine}{\textsc{Bourbaki Engine}}
\newcommand{\mirador}{\textsc{Mirador}}
\newcommand{\pc}{\textsc{ProofContext}}
\newcommand{\state}[1]{\texttt{[#1]}}
\newcommand{\otype}[1]{\textsc{#1}}
\newcommand{\op}[1]{\texttt{#1}}
\newcommand{\yes}{$\bullet$}
\newcommand{\pt}{$\circ$}
\newcommand{\no}{--}

\title[The Bourbaki Engine]{The Bourbaki Engine: A Governed Discovery
Architecture\\ for Long-Horizon Autonomous Mathematics}
\author{Maximiliano Lucius}
\address{Instituto Tecnol\'ogico de Buenos Aires, CABA, Argentina}
\email{maximiliano.lucius@tutamail.com}
\author{Jos\'e Minich}
\address{Instituto Tecnol\'ogico de Buenos Aires, C\'ordoba, Argentina}
\email{jminich@gmail.com}
\date{Draft --- \today}

\begin{document}

\begin{abstract}
Autonomous systems increasingly generate plausible research-level mathematics.
A prior architecture, the \engine{}, took the \emph{admission} problem as its
object: under what protocol should a generated claim be allowed to become a
load-bearing dependency of future work? That question is necessary but not
sufficient for \emph{discovery}: governing accumulation says nothing about how
new structure is produced over a research program that outlives any single
context. This paper reorganizes the engine around a dual thesis: reliable
autonomous mathematics requires two coupled but \emph{institutionally separated}
systems --- a \emph{Discovery Plane} that expands the space of mathematical
possibilities (and may be exploratory, redundant, speculative, and internally
contradictory) and a \emph{Trust Plane} that governs which outputs may become
reusable knowledge (and must remain conservative, auditable, and controlled by a
small explicit trusted computing base). We preserve and strengthen the original
contribution --- claim promotion as the unit of trusted progress, edge-level
soundness obligations, an epistemic separation of powers, fresh-context
independence as a falsifiable hypothesis, certificate discipline
(search~$\neq$~evidence~$\neq$~certificate~$\neq$~checker), an event-sourced
knowledge substrate, dependency-driven invalidation to a \state{SUSPECT} state,
and the False Theorem Promotion Rate as the primary trust metric. Around it we
specify seven modules with explicit interfaces (two stated as formal contracts),
two of which are implemented and separately published as companion preprints
(not yet peer-reviewed): \mirador{}, a typed, versioned representational
substrate whose unit is a Mathematical Cell (prototype and case study), and
\pc{}, a dependency-aware retrieval substrate returning verifiable Evidence
Bundles rather than similar text (prototype with controlled evaluation). The
remaining modules --- a Discovery Plane of machine-actionable operators, a
formalization plane, a certified computational laboratory, an enlarged
adversarial Congress, and a hierarchical research scheduler --- are specified at
explicitly labeled maturity levels. We replace the goal graph with a typed
research hypergraph in which object type is orthogonal to epistemic state, recast
the sixteen-round Bang affine plank campaign as a retrospective vertical slice
(requirements elicitation, not validation of unimplemented components), and
replace ambition-by-assertion with a gated readiness ladder (R0--R9) and a
three-track program. The engine does not promise to solve a Millennium Prize
Problem; a Millennium problem is a gated frontier stress test whose success
cannot be self-declared. The reusable asset is horizontal: a
verified-reasoning infrastructure applicable wherever claims are cheap to
generate and expensive to trust.
\end{abstract}

\maketitle

\section{Introduction}\label{sec:intro}

Mathematics has always scaled by dividing labour under a shared standard of
rigor. Hilbert's 1900 address turned open problems into public
milestones~\cite{Hilbert1900}; P\'olya systematized the heuristics of attack
\cite{Polya45}; Lakatos showed that proofs grow by surviving
refutations~\cite{Lakatos76}; and the Polymath experiments demonstrated that
loosely coordinated crowds can close serious theorems when the work is cut into
visible, checkable pieces~\cite{GowersNielsen09,Polymath12}. The most successful
sustained instance of industrialized rigor remains the Bourbaki collective:
many mathematicians, one identity, drafts destroyed and rewritten under
adversarial assembly until unanimity, and a cumulative corpus --- the
\emph{\'El\'ements} --- in which every statement rests on what came
before~\cite{Bourbaki,Mashaal06,Corry92,Weil92,Beaulieu93}.

Large-language-model agents now produce competent mathematical text, search
literature, write exact-arithmetic experiments, and drive proof
assistants~\cite{Yao23,Trinh24,RomeraParedes24,Davies21}. As generation
improves, two binding constraints appear, and they are different problems. The
first is \emph{admission}: deciding under what protocol a generated argument may
become a load-bearing dependency. The prior version of this
architecture~\cite{Lucius26} answered exactly that, and its answer stands: the
unit of progress is not a generated proof but a claim that survives a governed
promotion process. The second constraint is \emph{discovery}: producing new
mathematical structure --- definitions, reductions, barriers, conjectures,
proofs --- across a research program that spans many contexts and many months.
Governing accumulation, by itself, does not generate structure. A machine that
only prevents false promotions can sit forever at an empty frontier.

\begin{principle}[dual thesis]\label{pr:dual}
Reliable autonomous mathematics requires two coupled but institutionally
separated systems: a \emph{Discovery Plane} that expands the space of
mathematical possibilities, and a \emph{Trust Plane} that governs which outputs
may become reusable knowledge. The Discovery Plane may be exploratory,
redundant, speculative, and internally contradictory; the Trust Plane must
remain conservative, auditable, and controlled by a small explicit trusted
computing base. Neither may hold the other's authority.
\end{principle}

This reframes the platform without weakening its original spine. The Trust Plane
\emph{is} the claim-promotion architecture, unchanged in its guarantees and
still model-agnostic about how candidates are produced. What is new is that the
platform now also specifies the machinery of discovery --- representation,
retrieval, generation operators, formalization, and certified computation ---
and the interfaces by which their outputs enter the promotion boundary. Two of
those modules are not proposals but implemented, separately published companion
substrates: the \mirador{} representational layer~\cite{lucius2026mirador} and
the \pc{} retrieval layer~\cite{lucius2026proofcontext}. The rest are specified
at maturity levels we mark explicitly (\S\ref{sec:scope}), because the failure
this paper most wants to avoid is describing unbuilt machinery in a way that
reads as built.

The remaining sections state scope and contribution with maturity
(\S\ref{sec:scope}); locate the work against related systems by capability and
authority (\S\ref{sec:related}); specify the two-plane architecture and its
closed loop (\S\ref{sec:planes}); define the campaign constitution
(\S\ref{sec:constitution}), the \mirador{} substrate and the typed research
hypergraph (\S\ref{sec:mirador}), the \pc{} substrate and Evidence Bundles
(\S\ref{sec:proofcontext}), the Discovery Plane operators (\S\ref{sec:discovery}),
formalization and certified computation (\S\ref{sec:formal}), the enlarged
Congress (\S\ref{sec:congress}), and the hierarchical scheduler
(\S\ref{sec:scheduler}); persist all of it in an event-sourced substrate with
invalidation (\S\ref{sec:substrate}); encode the Bang pilot retrospectively
(\S\ref{sec:pilot}); replace the old validation plan with a readiness ladder and
three-track program (\S\ref{sec:validation}) and a gated Millennium stress test
(\S\ref{sec:millennium}); and state threats and limitations
(\S\ref{sec:threats}).

\section{Scientific scope and contribution}\label{sec:scope}

The revised paper answers a broader question than its predecessor:
\begin{quote}
How can autonomous systems generate, retrieve, test, formalize, adversarially
review, and cumulatively promote mathematical progress over long-horizon
research programs without allowing errors to become load-bearing knowledge?
\end{quote}
This is an architecture and research-program paper grounded in a
requirements-elicitation pilot, with falsifiable evaluation protocols and
explicitly separated implemented, manually exercised, specified, and aspirational
components. The contributions are:

\begin{enumerate}
\item[\textbf{(A)}] \textbf{The two-plane thesis} (\S\ref{sec:planes}):
discovery and trust as separated authorities, with a model-agnostic promotion
boundary between them.
\item[\textbf{(B)}] \textbf{Edge-level soundness obligations and a typed research
hypergraph} (\S\ref{sec:mirador}): a reduction or transformation
``$A$ follows from $B_1,\dots,B_k$'' is itself an auditable claim, and object
type (what a node \emph{is}) is orthogonal to epistemic state (how far it has
\emph{survived}).
\item[\textbf{(C)}] \textbf{Two implemented substrates with bounded contracts}:
\mirador{}~\cite{lucius2026mirador} for representation and
\pc{}~\cite{lucius2026proofcontext} for retrieval, stated as interfaces rather than
reproduced in full.
\item[\textbf{(D)}] \textbf{A Discovery Plane of machine-actionable operators}
(\S\ref{sec:discovery}), each with preconditions, produced cell type, initial
state, required retrieval context, validation obligation, and predictable
failure mode.
\item[\textbf{(E)}] \textbf{An epistemic separation of powers and an enlarged
adversarial Congress} (\S\ref{sec:congress}): generation, promotion, and
certification are never vested in one agent; fresh-context independence is a
falsifiable hypothesis, not an intuition.
\item[\textbf{(F)}] \textbf{Certificate discipline as a distinct trust layer}
(\S\ref{sec:formal}): search~$\neq$~evidence~$\neq$~certificate~$\neq$~checker,
with a small explicit trusted computing base and an independent
statement--formalization correspondence audit.
\item[\textbf{(G)}] \textbf{A gated evaluation program} (\S\ref{sec:validation},
\S\ref{sec:millennium}): the primary metric is the False Theorem Promotion Rate;
progress toward frontier problems is gated by a readiness ladder that the engine
cannot self-declare complete.
\end{enumerate}

\noindent The discipline these contributions demand of \emph{us} is a maturity
ledger. Every mechanism occupies exactly the column its evidence supports;
future-tense architecture is never written as if implemented.

{\footnotesize
\begin{longtable}{@{}p{0.30\textwidth}ccccc@{}}
\caption{Maturity ledger. C = conceptual/specified; M = manually exercised
(human-orchestrated pilot); P = prototype implemented; E = controlled
evaluation; X = external validation. The highest supported level is marked.}
\label{tab:maturity}\\
\toprule
\textbf{Mechanism} & \textbf{C} & \textbf{M} & \textbf{P} & \textbf{E} &
\textbf{X}\\
\midrule
\endfirsthead
\toprule
\textbf{Mechanism} & \textbf{C} & \textbf{M} & \textbf{P} & \textbf{E} &
\textbf{X}\\ \midrule
\endhead
\mirador{} representation~\cite{lucius2026mirador} & & & \yes & & \\
\pc{} retrieval~\cite{lucius2026proofcontext} & & & & \yes & \\
Event log, hash-chaining, invalidation (in \mirador{}) & & & \yes & & \\
Claim promotion / separation of powers & & \yes & & & \\
Adversarial Congress (fresh-context review) & & \yes & & & \\
Certified computation (one instance: $812{,}651$-leaf certificate + checker) & &
\yes & & & \\
Discovery Plane operators & \yes & & & & \\
Formalization plane (Lean/\texttt{mathlib} backend) & \yes & & & & \\
Hierarchical research scheduler & \yes & & & & \\
Campaign constitution & \yes & & & & \\
Millennium-scale campaign & \yes & & & & \\
\bottomrule
\end{longtable}
}

\noindent No row is marked \textbf{X}: nothing here has yet survived external
peer review as a system, and the Bang manuscript is not yet journal-accepted.
The engine cannot promote itself to that column (\S\ref{sec:validation}).

\section{Related work}\label{sec:related}

We locate the engine not by capability alone but by \emph{who holds which
authority}. Systems that ``do mathematics'' can be separated by whether they
generate, retrieve, formalize, certify, persist across campaigns, govern
promotion, invalidate dependencies, and build theory. Table~\ref{tab:matrix}
records this; every 2026 entry was verified against its primary source (access date 2026-07-12; the coding table is
supplementary material), and no capability coding is inferred
from an abstract. In particular the certification column separates
kernel-certified provers from LLM-verifier self-checking, which is not
certification.

{\scriptsize
\begin{longtable}{@{}p{0.24\textwidth}cccccccc@{}}
\caption{Capability-and-authority matrix. \yes~= provided; \pt~= partial or
scoped; \no~= not provided. Columns: \textbf{Gen}erate, \textbf{Retr}ieve,
\textbf{Form}alize, \textbf{Cert}ify, \textbf{Pers}ist across campaigns,
\textbf{Gov}ern promotion, \textbf{Inv}alidate dependencies, build
\textbf{Theory}. For the \engine{} and its substrates, a cell denotes capability
the architecture \emph{specifies} (cf.\ the maturity ledger,
Table~\ref{tab:maturity}); a \yes{} appears only where a substrate implements it
--- retrieval via \pc{}, invalidation via \mirador{} --- and manual or conceptual
capabilities are marked \pt{} at most, never at parity with a shipping system.}
\label{tab:matrix}\\
\toprule
\textbf{System} & Gen & Retr & Form & Cert & Pers & Gov & Inv & Theory\\
\midrule
\endfirsthead
\toprule
\textbf{System} & Gen & Retr & Form & Cert & Pers & Gov & Inv & Theory\\ \midrule
\endhead
Aletheia~\cite{Aletheia26} & \yes & \yes & \no & \no & \no & \no & \no & \no\\
QED~\cite{QED26} & \yes & \pt & \no & \no & \no & \pt & \pt & \no\\
RMA~\cite{RMA26} & \yes & \yes & \no & \no & \pt & \pt & \no & \no\\
Prover Agent~\cite{ProverAgent25} & \yes & \no & \yes & \yes & \no & \no & \no & \no\\
LEAP~\cite{LEAP26} & \yes & \no & \yes & \yes & \no & \no & \no & \no\\
FunSearch~\cite{RomeraParedes24} & \yes & \no & \no & \pt & \pt & \pt & \no & \no\\
AlphaEvolve~\cite{AlphaEvolve25} & \yes & \no & \no & \pt & \yes & \yes & \no & \no\\
Lean \texttt{mathlib}~\cite{mathlib20} & \no & \pt & \yes & \yes & \yes & \pt & \pt & \yes\\
iCoq~\cite{Celik17} & \no & \no & \yes & \yes & \yes & \no & \yes & \no\\
\texttt{leanblueprint}~\cite{Massot,Tao23} & \no & \no & \pt & \pt & \yes & \no & \pt & \pt\\
LeanSearch\,v2~\cite{LeanSearch26} & \pt & \yes & \no & \no & \no & \no & \no & \no\\
TheoremGraph~\cite{TheoremGraph26} & \pt & \yes & \no & \no & \pt & \no & \no & \pt\\
MaRDI~\cite{MaRDI23} & \no & \yes & \no & \no & \yes & \pt & \no & \pt\\
Zep/Graphiti~\cite{Zep25} & \no & \yes & \no & \no & \yes & \no & \pt & \no\\
\midrule
\mirador{}~\cite{lucius2026mirador} & \no & \pt & \pt & \no & \yes & \no & \yes & \no\\
\pc{}~\cite{lucius2026proofcontext} & \no & \yes & \no & \no & \pt & \no & \yes & \no\\
\engine{} (this paper) & \pt & \yes & \pt & \pt & \pt & \pt & \yes & \pt\\
\bottomrule
\end{longtable}
}

\noindent Three points survive this matrix. First, the recent autonomous
systems~\cite{Aletheia26,QED26,RMA26} primarily optimize the probability of
\emph{finding} a correct result; none maintains a governed cross-campaign
corpus. (Notably, QED enforces that its prover and verifiers share no context ---
an independent adoption of the separation this paper argues for.) Second, the
formal-methods lineage --- \texttt{mathlib}~\cite{mathlib20},
iCoq~\cite{Celik17}, \texttt{leanblueprint}~\cite{Massot,Tao23} --- supplies
persistence, certification, and dependency tracking, but inside a single
always-consistent type-checked environment that a corpus of possibly-refuted,
semi-formal claims cannot assume. Third, the knowledge-graph and agent-memory
systems~\cite{TheoremGraph26,MaRDI23,Zep25,Yadav26} persist and retrieve
structure but hold \emph{no} epistemic authority and do not govern promotion.
The engine's delta is the composition: it governs promotion \emph{and} builds
on a representation (\mirador{}) and retrieval (\pc{}) that carry epistemic
state and invalidate on it. Coding each entry from its primary source confirms
this gap is systemic: of the systems in Table~\ref{tab:matrix}, none carries a live
epistemic state (\state{OPEN}/\state{PROMOTED}/\state{REFUTED}) \emph{together
with} retrieval that hard-rejects an incompatible-but-similar match. Machine
proof search --- from the Logic Theory
Machine~\cite{NewellSimon57} and AND/OR search~\cite{Nilsson71} through proof
planning~\cite{Bundy88} to McCune's Robbins solution~\cite{McCune97} --- supplies
the AND--OR skeleton the hypergraph generalizes; the computer-assisted
lineage~\cite{AppelHaken77,Gonthier08,Hales17,Heule16,deBruijn70} fixes the
certificate standard of \S\ref{sec:formal}.

\section{The two-plane architecture}\label{sec:planes}

The engine is defined by a boundary, not by its agents. Candidate mathematical
objects are produced on one side and admitted, if at all, on the other. The
original single boundary --- ``candidate $\to$ evidence $\to$ promotion $\to$
corpus'' --- is now embedded in a closed loop:
\[
\begin{gathered}
\underbrace{\text{corpus}}_{\text{\'El\'ements}}
\to
\underbrace{\pc{}}_{\text{retrieval}}
\to
\underbrace{\text{Discovery Plane}}_{\text{operators}}
\to
\text{candidates}\\
\to
\underbrace{\text{formalize}+\text{certify}}_{}
\to
\boxed{\text{Trust Plane / Congress}}
\to
\text{corpus}\to\cdots
\end{gathered}
\]
Everything up to the box is Discovery; the box and the corpus behind it are
Trust. The membrane between them is one-way in authority: Discovery and
retrieval may \emph{propose} nodes and edges; only the Trust Plane may
\emph{promote}.

\begin{principle}[governance, not generation, at the boundary]\label{pr:governs}
The Trust Plane does not define how claims must be generated. It defines the
conditions under which claims may become load-bearing dependencies. Any
generator --- an end-to-end research system, a role-specialized solver, a single
proof model --- is a candidate source with respect to the boundary.
\end{principle}

What changes from the predecessor is that ``the engine does not define how claims
are generated'' is now true only \emph{of the Trust Plane}, not of the whole
platform. The integrated engine specifies the interfaces and governance of
discovery --- what a candidate is, where it is retrieved from, which operators
produce it, how it is formalized and certified --- without mandating a single
generation algorithm. The two planes have opposite virtues by design: the
Discovery Plane is judged by yield and diversity and is allowed to be wrong; the
Trust Plane is judged by its False Theorem Promotion Rate and is allowed to be
slow. Collapsing them --- letting a generator certify its own output --- is the
precise failure the architecture exists to prevent.

\section{The campaign constitution}\label{sec:constitution}

Before any graph expansion, a campaign is \emph{constituted}: a frozen,
hash-pinned document that fixes the rules of the game so that neither a
generating agent nor an over-eager scheduler can quietly move them after seeing
results. The constitution freezes: the exact target and admissible variants;
definitions and conventions; assumptions and excluded regimes; \emph{what counts
as a solution}; allowed foundational backends and axioms; certification
requirements by claim type; the novelty and attribution policy; evaluation
metrics; compute and human-review budget; stop, retreat, and publication
conditions; and the minimum useful deliverables.

\begin{principle}[solution is defined outside the engine]\label{pr:solution}
For a hard target --- above all a Millennium problem --- ``solution'' cannot be
defined by an internal agent verdict. The constitution must name the external
process (journal acceptance, recognized community validation) that alone can
close the top of the readiness ladder (\S\ref{sec:validation}).
\end{principle}

The constitution is the object that makes the pilot's costliest lesson
un-repeatable by construction: a target stated as the unrestricted constant $C$
but attacked for the one-plank-per-direction constant $C^{111}$ is a
constitution violation detectable before a single proof is written
(\S\ref{sec:pilot}).

\section{\mirador{}: representation and the typed research hypergraph}\label{sec:mirador}

The substrate over which everything is stored is \mirador{}, a typed, versioned
Mathematical Intermediate Representation whose unit is a \emph{Mathematical
Cell}~\cite{lucius2026mirador}. The engine depends only on its interface, not on
its internals; \mirador{} is a separately published prototype with an executable
schema, a dependency-free reference implementation, twenty-eight deterministic
invariant tests, and a thirteen-cell case study on the Bang problem.

\begin{contract}[\mirador{} interface]\label{con:mirador}
The representation must provide: \emph{Mathematical Cells} binding a canonical
claim to its explicit context; \emph{object type orthogonal to epistemic state};
\emph{typed dependency and transformation edges} as first-class objects, each
carrying its own soundness obligation; \emph{versioning with semantic diffs};
claim identity as the canonical hash of a statement \emph{and its context}
(assumptions and scope participate in identity); and links between natural
language, computation, and proof-assistant declarations that are recorded, never
inferred from compilation.
\end{contract}

The engine's goal graph is generalized, in the design, to a \emph{typed research
hypergraph} over these cells. Two orthogonal axes classify every node.

\paragraph{Object type (what a node is).}
\otype{DEFINITION}, \otype{QUESTION}, \otype{CONJECTURE}, \otype{CLAIM},
\otype{REDUCTION}, \otype{EQUIVALENCE}, \otype{EXAMPLE}, \otype{COUNTEREXAMPLE},
\otype{BARRIER}, \otype{EXPERIMENT}, \otype{CERTIFICATE}, \otype{CHECKER},
\otype{HEURISTIC}, \otype{TRANSLATION}, \otype{PROGRAM}.

\paragraph{Epistemic state (how far it has survived).}
\state{DRAFT}, \state{OPEN}, \state{EVIDENCE}, \state{CANDIDATE},
\state{REFUTED}, \state{NO\_GO}, \state{FORMALIZED}, \state{AUDITED},
\state{PROMOTED}, \state{SUSPECT}, \state{RETIRED}.

\begin{principle}[type is not status]\label{pr:orthogonal}
An object typed \otype{CLAIM} and labelled ``Theorem'' in prose is not thereby
\state{PROMOTED}. The literary label names the \emph{intent}; the epistemic
state names the \emph{evidence}. Conflating them is the mechanism by which a
plausible draft is cited as a theorem, and the hypergraph forbids it by keeping
the axes separate.
\end{principle}

The inverse AND--OR proof structure of the predecessor is retained as one
important \emph{subgraph}: a \otype{REDUCTION} node with an AND-set of premises
is exactly an AND-node, and alternative reductions to the same goal are an
OR-node. But the hypergraph also represents questions, definitions, experiments,
barriers, equivalences, translations, and entire research programs that the pure
proof tree could not. Every \otype{REDUCTION}, \otype{EQUIVALENCE}, and
\otype{TRANSLATION} edge carries an auditable soundness obligation, checked
separately from the nodes it joins --- the boundary at which the pilot's most
expensive defect actually lived (\S\ref{sec:pilot}). Two standing graph audits
are specified: \emph{acyclicity of justification} (no node reachable from its own proof
obligations) and \emph{definition--proof consistency} (the object named in a
statement is byte-identical, by hash, with the object manipulated in its proof).

\section{\pc{}: dependency-aware retrieval and Evidence Bundles}\label{sec:proofcontext}

Discovery and proof both begin by asking the corpus what is already known. In
mathematics, similarity is unsafe: two statements can be nearly identical in
wording yet differ in a quantifier, a dimension, a symmetry assumption, a
normalization, a definition version, or an epistemic status. \pc{} is a
dependency-aware retriever, layered on \mirador{}, that returns not the most
similar passage but the smallest context in which a claim can be
\emph{safely used}~\cite{lucius2026proofcontext}. It is a separately published
prototype with a controlled evaluation.

\begin{contract}[\pc{} interface]\label{con:pc}
Given a \emph{Problem Signature} compiled from a natural-language query,
retrieval returns a typed, hashed \emph{Evidence Bundle}: definitions,
assumptions, and exclusions; relevant claims at exact versions; dependency
paths; proofs, certificates, and formal declarations; counterexamples and
barriers; a provenance manifest; and epistemic states with warnings. Retrieval
may propose candidate cells and edges and may rank and annotate them; it holds
\emph{no promotion authority} (the authority membrane). Filters that reject a
candidate incompatible on a typed axis do so regardless of its similarity score,
and both indexes and emitted bundles are invalidated version-aware.
\end{contract}

On \textsc{MathScopeBench}, an adversarial benchmark of mathematical near-misses,
\pc{} serves $0\%$ of incompatible near-matches against $44$--$75\%$ for lexical,
dense, hybrid, and hybrid-without-filters retrieval, at equal or higher recall;
the advantage disappears exactly when the hard filters are removed, isolating
them as the cause. On downstream proxies, dependency sufficiency rises to
$0.90$--$0.93$ (from $0.40$--$0.43$ for a flat hybrid) and correct refusal to
$1.00$ (from $0.00$). These are prototype-scale numbers on a self-built
benchmark whose near-miss categories were authored alongside the very filters
they reward: the ablation isolates causal attribution but not this
construct-validity risk, so a held-out or third-party benchmark is required
before R1 is treated as robust. We mark them as such (Table~\ref{tab:maturity});
they establish retrieval readiness (R1, \S\ref{sec:validation}), not the engine.

\section{The Discovery Plane}\label{sec:discovery}

The Discovery Plane is a library of machine-actionable \emph{operators} that
take retrieved Evidence Bundles and emit new \mirador{} cells in an initial
epistemic state. Operators are the specified, not-yet-implemented heart of the
generative side (Table~\ref{tab:maturity}); the contribution here is the
\emph{contract}, which forces every operator to declare what would make its
output wrong.

{\footnotesize
\begin{longtable}{@{}p{0.20\textwidth}p{0.30\textwidth}p{0.40\textwidth}@{}}
\caption{Discovery operators (excerpt). Each declares admissible input,
produced cell type and initial state, and the validation obligation that its
output must later discharge at the Congress. All emit in \state{DRAFT} or
\state{CANDIDATE}; none may emit \state{PROMOTED}. The full six-field per-operator
contract (adding precondition and required retrieval context) is in the
supplementary specification.}\label{tab:ops}\\
\toprule
\textbf{Operator} & \textbf{Produces} & \textbf{Validation obligation /
predictable failure}\\ \midrule
\endfirsthead
\toprule
\textbf{Operator} & \textbf{Produces} & \textbf{Validation obligation /
predictable failure}\\ \midrule
\endhead
\op{SPECIALIZE} & a \otype{CLAIM}/\otype{EXAMPLE} in a restricted regime &
soundness of the restriction edge; failure: a true special case marketed as the
general one.\\
\op{GENERALIZE} & a \otype{CONJECTURE} over a wider class & a counterexample
search must be scheduled; failure: over-generalization past a hidden barrier.\\
\op{COUNTEREXAMPLE\_FIRST} & a \otype{COUNTEREXAMPLE} or its absence as
\state{EVIDENCE} & exact-arithmetic reproduction; failure: numeric near-miss
mistaken for a counterexample.\\
\op{NEGATE\_AND\_SEARCH} & a \otype{BARRIER} node (state \state{NO\_GO}) &
the impossibility must be proved, not merely unfound; failure: absence of
evidence read as evidence of absence.\\
\op{DUALIZE} & an \otype{EQUIVALENCE} to a dual formulation & the duality edge's
soundness; failure: a lossy dual asserted as exact.\\
\op{CHANGE\_REPRESENTATION} & a \otype{TRANSLATION} cell & statement--translation
correspondence audit; failure: a weaker statement encoded than intended.\\
\op{TRANSFER} & a \otype{CLAIM} imported by analogy from another domain &
the analogy is a \state{DRAFT} until its edge is discharged; failure: false
transfer across a disanalogy.\\
\op{MINIMAL\_MISSING\_LEMMA} & a \otype{CLAIM} whose proof would close a goal &
acyclicity audit; failure: assuming the conclusion.\\
\op{EXTREMALIZE} / \op{STABILITY} & an \otype{EXPERIMENT} probing extremal or
perturbed cases & labelled \state{EVIDENCE}, never proof; failure: a grid
mistaken for a theorem.\\
\op{BARRIER\_SEARCH} & a \otype{BARRIER} capping a method family & sharpness must
be proved; failure: a soft obstruction sold as a hard one.\\
\op{FORMALIZE\_EARLY} & a \otype{TRANSLATION} to a proof assistant & kernel check
plus correspondence audit; failure: green kernel, wrong theorem.\\
\op{SYNTHESIZE\_DEFINITION} & a \otype{DEFINITION} cell (new version) & triggers
definition--proof re-audit of dependents; failure: definition drift.\\
\bottomrule
\end{longtable}
}

\begin{principle}[candidate generation is not novelty]\label{pr:novelty}
An operator produces a \emph{candidate}, not a discovery. Genuine mathematical
novelty is adjudicated downstream, by a novelty-and-priority audit against the
literature (\S\ref{sec:congress}) and by promotion. Superficial novelty by
renaming is a threat the architecture must detect (\S\ref{sec:threats}), not a
success the Discovery Plane may claim.
\end{principle}

Two loops drive the operators. The \emph{deductive loop} takes an \state{OPEN}
claim, retrieves its Evidence Bundle, proposes a reduction or proof plan,
generates obligations, searches for counterexamples, formalizes or proves,
submits to the Congress, and receives promotion or a repair order. The
\emph{experimental loop} takes a mathematical family, generates exact instances,
computes or searches, detects structure, synthesizes a conjecture or definition,
attacks it with counterexamples, and extracts a certificate or a proof
obligation before the same Congress. Every product of either loop enters the
event log as typed events and cells.

\section{Formalization and certified computation}\label{sec:formal}

On the verification axis the engine claims no novelty: it places itself in the
Flyspeck / SAT-proof / Lean-\texttt{mathlib} lineage of computer-assisted
mathematics~\cite{Hales17,Heule16,Gonthier08,mathlib20} and treats that
lineage's rules as non-negotiable. What it adds is to make verification a
distinct trust layer with a small explicit trusted computing base (TCB), and to
audit two correspondences the lineage usually leaves implicit.

\paragraph{Certified computation.}
The admission policy requires computation to enter the corpus only in Flyspeck
form: the machine-dependent step
is isolated in a single finite lemma with a \emph{static} certificate (e.g.\ a
branch-and-bound tree whose every leaf is discharged by a sound predicate),
verified by an \emph{independent} checker (no shared code with the search,
standard library only, exact rational arithmetic), with mutation tests
establishing the checker's non-vacuity and full hashes shipped with the artifact.
\begin{principle}[discard the discoverer]\label{pr:discard}
Trust in a computational result rests on the certificate and its independent
checker alone. The search, its heuristics, and its transcript may be discarded
in full without affecting verifiability:
search~$\neq$~evidence~$\neq$~certificate~$\neq$~checker. Every laboratory
output is labelled \emph{heuristic}, \emph{evidence}, \emph{certificate}, or
\emph{proof}, and the label is load-bearing.
\end{principle}
The pilot's flagship instance is an $812{,}651$-leaf exact-rational certificate
for a sharp one-plank-per-direction covering constant, verified by a one-page
independent checker; four hand-proved soundness lemmas and one machine-checked
emptiness fact are the human/machine boundary~\cite{Lucius26}. A cautionary
counterweight is on record: Ambrus's withdrawn de-computerization of Ball's plank
theorem~\cite{Ambrus22} reminds the engine that removing a certificate fails at
a measurable rate, which is why certificates are never replaced by prose without
a full Congress cycle.

\paragraph{Formalization.}
Formalization is specified to run as a continuous pipeline: informal statement $\to$ \mirador{}
IR $\to$ formal declaration $\to$ proof state $\to$ kernel feedback $\to$
localized repair. Lean~4/\texttt{mathlib}~\cite{mathlib20} is the initial
backend, but the architecture is backend-aware, not backend-identical.
\begin{principle}[a kernel check is not a correspondence proof]\label{pr:correspond}
That a formal declaration type-checks does \emph{not} establish that it encodes
the intended informal theorem. Statement--formalization correspondence is an
independent claim, audited by a dedicated Congress lens (\S\ref{sec:congress})
and recorded in \mirador{} as \texttt{asserted}, never inferred from
compilation.
\end{principle}

\section{The Congress: roles, powers, and adversarial lenses}\label{sec:congress}

Agents are cheap; the architecture spends them on opposition. The original roles
are preserved: \textbf{Strat\`ege} (expands the graph, maintains the
\emph{plan de travail}), \textbf{Scribe} (attacks an \state{OPEN} node and
drafts the proof), \textbf{Exp\'erimentateur} (exact-arithmetic experiments; a
numeric grid is never a proof), \textbf{Biblioth\'ecaire} (deep literature
research with adversarial source verification), \textbf{R\'efutateurs} (the
panel proper), \textbf{Auditeur} (standing graph-level auditor), and
\textbf{Greffier} (owns certificates, checkers, and the append-only record).

\begin{principle}[epistemic separation of powers]\label{pr:sep}
Generation, promotion, and certification are distinct authorities and are never
vested in the same agent. The authority that generates a claim
(\textbf{i})~does not promote it, (\textbf{ii})~does not choose its verdict,
(\textbf{iii})~does not control the checker that verifies its computation, and
(\textbf{iv})~does not contaminate the referee with the search transcript. This
constrains \emph{who may exercise which power}, not merely how many agents
participate.
\end{principle}

\begin{principle}[unanimity at the boundary]\label{pr:veto}
Promotion to \state{PROMOTED} requires the Congress: any panelist can veto, and
the veto is a \emph{structured object} --- target node/edge, concrete objection,
affected dependency paths, and discharge condition --- returned to the Scribe as
a work order, since the Scribe holds no promotion authority over its own node.
\end{principle}

The enlarged frontier of the two-plane engine demands new lenses beyond local
correctness. The specified panel adds, as distinct fresh-context briefs not yet
exercised (the pilot predates them): an
\emph{informal--formal correspondence} auditor (Principle~\ref{pr:correspond});
a \emph{novelty and priority} auditor (Principle~\ref{pr:novelty}); a
\emph{retrieval-provenance} auditor (does the Evidence Bundle actually support
what was drawn from it?); an \emph{assumption/scope} auditor (the $C$ versus
$C^{111}$ lens); a \emph{checker-independence} auditor (no common-mode code
between search and checker); and an \emph{integration and invalidation} auditor
(does promotion trigger the correct downstream re-audits?).

\begin{principle}[fresh-context independence, as a hypothesis]\label{pr:fresh}
Write $P(\text{defect detected}\mid c)$ for the probability a review detects a
genuine defect given the reviewer's epistemic context $c$. The engine's wager is
$P(\text{detected}\mid\text{fresh}) > P(\text{detected}\mid\text{shared
generation history})$, along a ladder of increasing independence (self-review;
watched; single fresh-context; multi-lens panel). Whether climbing this ladder
measurably improves verification is an empirical question the pilot motivates
($n=1$) but cannot settle, and that \S\ref{sec:validation} proposes to measure.
\end{principle}

\section{Hierarchical research scheduling}\label{sec:scheduler}

A flat frontier bandit over open nodes cannot run a years-long program. The
scheduler is specified to be hierarchical, allocating a portfolio at three levels:
(1)~\emph{tactical} proof obligations; (2)~\emph{research programs} lasting many
cycles; and (3)~\emph{meta-program} decisions about representations and entire
method families. At each level it balances exploitation, orthogonal exploration,
counterexample search, formalization, certified computation, literature and
priority research, \state{NO\_GO} attempts, reproduction, and minimum publishable
milestones. Within a program, the classical index
\[
\text{priority}(v)=\frac{\widehat{P}(\text{close }v)\cdot\text{value}(v\to R)}
{\text{cost}(v)}
\]
would rank tactical nodes, with $\widehat{P}$ updated in bandit
fashion~\cite{Auer02,Browne12} and value flowing down from the root milestone
$R$ through the AND--OR substructure.

\begin{principle}[the score is not evidence]\label{pr:score}
A scheduler priority is a resource-allocation heuristic, never evidence of truth.
It may decide what to attempt next; it may never decide what is promoted.
\end{principle}

Three hard rules temper the index, unchanged from the predecessor: \emph{time
boxing} (an exhausted attempt returns to the Strat\`ege with its transcript,
never to the same Scribe for ``one more try''); the \emph{retreat rule} (after
$N$ failures a node is re-entered only via a new reduction edge or a
\state{NO\_GO} attempt --- a proved \otype{BARRIER} (a route reaching state
\state{NO\_GO}) being a first-class asset that prunes infinite search); and
\emph{milestone insurance} (the graph always contains a resourced path to a
shippable \emph{fascicule}, so a campaign degrades gracefully).

\section{Event-sourced knowledge and invalidation}\label{sec:substrate}

A governance architecture is only as trustworthy as the record it governs. The
substrate does not persist knowledge as prose.

\begin{principle}[what memory holds]\label{pr:memory}
The engine remembers \emph{epistemic state}, \emph{dependency},
\emph{provenance}, and \emph{confidence} --- where confidence is the pair
(state, review level). Prose and transcripts are attachments to that record,
never the record itself.
\end{principle}

Four storage roles realize this, with a strict direction of authority; the core
of the first three is what \mirador{} implements.
\emph{(i)~The log is the source of truth.} An append-only sequence of campaign
events (\texttt{proposed}, \texttt{evidence-added}, \texttt{promoted} recording
review-independence level, \texttt{refuted}, \texttt{obligation-discharged},
\texttt{certificate-registered}, \texttt{retired}); every other store is a
derived, reconstructible view~\cite{Fowler05}, made tamper-evident by
hash-chaining with periodic Merkle roots, the transparent-log pattern of
certificate transparency~\cite{RFC6962,Cox19}. Rebuildability buys time travel
and auditability of the promotion process itself --- the metrics of
\S\ref{sec:validation} are computable precisely because each promotion event
carries its review level.
\emph{(ii)~The graph is a truth-maintained cache.} Typed edges, each carrying its
soundness obligation; its defining operation is retraction propagation: when a
node is refuted or a definition it rests on is revised, invalidation propagates
through the upward transitive closure, demoting every affected \state{PROMOTED}
node to \state{SUSPECT} with re-audit work orders. (The propagation primitive is
implemented in \mirador{}; the end-to-end promotion loop that would act on the
work orders is not.) \state{SUSPECT} is the only state entered automatically; the
only exit is back through the Congress. This is
the truth-maintenance tradition~\cite{Doyle79,deKleer86}; the closest working
precedent is iCoq~\cite{Celik17}, which invalidates the upward closure of changed
entities. \texttt{leanblueprint}~\cite{Massot,Tao23} propagates \emph{good} news
up a dependency graph; the engine specifies the complementary propagation of
\emph{bad} news as first-class epistemic demotion, distinct from iCoq's
recheck-invalidation, which \texttt{mathlib} can forgo only because it lives in
one always-consistent environment~\cite{vanDoorn20}.
\emph{(iii)~The corpus is immutable and content-addressed.} \emph{Les
\'El\'ements} holds only promoted objects, frozen at promotion; identity is the
canonical hash of the statement-with-context (the \mirador{} identity rule),
so that ``object named in the statement $\ne$ object manipulated in the proof''
becomes a comparison of two hashes, and citations pin hashes in the manner of
trusty URIs~\cite{Kuhn14,Kuhn16}. Provenance reuses PROV-DM
vocabulary~\cite{PROVDM13}.
\emph{(iv)~Heuristic memory has no authority.} Embeddings and traces are an index
\emph{into} the graph; retrieval must resolve to a node and its epistemic state
before use. The membrane is one-way: soft memory may propose an \state{OPEN}
node, never assert a \state{PROMOTED} one --- the failure it prevents is a
hallucinated recollection leaking into a proof. Recent agent-memory
systems~\cite{Zep25,Yadav26} manage validity \emph{in time} (``no longer
holds''), not epistemic refutation (``proved false''), and separate no authority
between soft memory and a verified corpus; for the engine that separation is a
requirement.

The one binding physical requirement is transactional: a promotion commits its
event, the state change, the corpus insertion, and the certificate registration
atomically, or not at all. The separation of powers (Principle~\ref{pr:sep})
is specified to materialize as database permissions --- the generating role
would hold no write access to the corpus.

\section{The Bang affine plank vertical slice}\label{sec:pilot}

The pilot is a \emph{requirements-elicitation case study}, not an evaluation, and
its encoding here is \emph{retrospective}: \mirador{}, \pc{}, and the Discovery
Plane operators were \emph{not} implemented during the 2026 sixteen-round manual
campaign on Bang's affine plank conjecture~\cite{Bang51,Ball91,BY26,Verreault26}.
We show how the revised engine \emph{would} represent that campaign's documented
events; we do not claim it ran them. With $n=1$ the campaign motivates the
mechanisms and cannot measure them.

The campaign's outputs were a 35-page manuscript~\cite{Lucius26} (a
transport-defect framework with minimax duality; a characterization of witness
measures for three directions on the triangle, occupying the gap left by
Gardner~\cite{Gardner88}; a rigid two-parameter family of sharp coverings; a
dimension-general facet defect; and the certified covering constant of
\S\ref{sec:formal}), a verified state-of-the-art map~\cite{BY26,AKP19,GKP22,MOM26},
and a ranked program of next attacks. The design-relevant data are the
\emph{failures}, all caught before publication, none inside a proof, and each now
owned by a specific mechanism and a specific typed transition.

{\footnotesize
\begin{longtable}{@{}p{0.40\textwidth}p{0.24\textwidth}p{0.28\textwidth}@{}}
\caption{Retrospective encoding of pilot events --- documented defects and key
artifacts. Each maps to an object/edge in the typed hypergraph, an epistemic
transition, and the owning mechanism.}\label{tab:pilot}\\
\toprule
\textbf{Pilot event (defect or key artifact)} & \textbf{Hypergraph
locus} & \textbf{Mechanism / transition}\\ \midrule
\endfirsthead
\toprule
\textbf{Pilot event} & \textbf{Hypergraph locus} & \textbf{Mechanism}\\ \midrule
\endhead
Object proved ($C^{111}$) $\ne$ object defined ($C$) & two \otype{DEFINITION}
cells, distinct context-hashes & assumption/scope lens; constitution violation
(\S\ref{sec:constitution})\\
Invalid reduction edge (proved for $C^{111}$, stated for $C$) &
\otype{REDUCTION} edge obligation & edge-level audit; $\to$\state{REFUTED} edge
(\S\ref{sec:mirador})\\
Circular termination argument & \otype{CLAIM} reachable from its own obligations
& acyclicity audit; static-certificate template (\S\ref{sec:formal})\\
Priority claim (``first non-concurrent triple'') falsified & \otype{CLAIM} vs
\otype{COUNTEREXAMPLE} & novelty/priority lens (\S\ref{sec:congress})\\
Six internal audits missed what one external pass caught & review-level metadata
on the \texttt{promoted} events & fresh-context panels; \emph{la retraite}
(Principle~\ref{pr:fresh})\\
Sign-cell typo in a soundness lemma; degenerate-direction gap &
\otype{CLAIM}/\otype{CHECKER} with a \dbend{} flag & multi-lens re-derivation;
dangerous-bend discharge\\
$812{,}651$-leaf certificate + independent checker & \otype{CERTIFICATE},
\otype{CHECKER} cells & certificate discipline (Principle~\ref{pr:discard})\\
Route obstructions (Bang sign argument; symplectic ceiling) & \otype{BARRIER}
nodes (state \state{NO\_GO}) & retreat rule; barrier as asset
(\S\ref{sec:scheduler})\\
\bottomrule
\end{longtable}
}

Conversely, what never failed --- $\sim$22 adversarially re-derived computations,
massive exact-arithmetic cross-checks, the certificate/checker pair --- is
exactly what had already been subjected to opposition and exact arithmetic. The
working hypothesis --- \emph{errors concentrate where claims meet language, so
that is where adversaries should be spent} --- is consistent with the one
campaign on record but remains a hypothesis until the node-versus-edge error
share of \S\ref{sec:validation} is measured across many.

\section{Validation ladder and three-track program}\label{sec:validation}

We replace any impression that the next step is ``solve a Millennium problem''
with a gated readiness ladder. Each rung is a falsifiable claim about the system,
and the engine cannot self-declare the top rungs.

{\footnotesize
\begin{longtable}{@{}p{0.07\textwidth}p{0.85\textwidth}@{}}
\caption{The readiness ladder. Current status is marked against
Table~\ref{tab:maturity}: R0--R1 have prototype evidence; R2--R9 are
open.}\label{tab:ladder}\\
\toprule
\textbf{Rung} & \textbf{Readiness criterion}\\ \midrule
\endfirsthead
\toprule \textbf{Rung} & \textbf{Readiness criterion}\\ \midrule \endhead
R0 & \emph{Representation}: \mirador{} schemas, identity, and versioning operate
on real claims. \textbf{(prototype)}\\
R1 & \emph{Retrieval}: \pc{} rejects incompatible near-matches and reconstructs
usable Evidence Bundles. \textbf{(prototype eval)}\\
R2 & \emph{Governance}: events, permissions, promotion, veto, invalidation, and
deterministic replay work end to end.\\
R3 & \emph{Discovery}: the system regenerates nontrivial known results under
contamination controls and records failed routes.\\
R4 & \emph{Open-problem}: a genuinely new lemma, counterexample, reduction, or
barrier on a bounded open problem.\\
R5 & \emph{External}: at least one result survives independent expert or journal
review.\\
R6 & \emph{Millennium mapping}: a target-specific constitution and a verified
state-of-the-art atlas are released.\\
R7 & \emph{Millennium frontier}: a novel, externally reviewed result advances a
recognized subproblem or proves a useful barrier.\\
R8 & \emph{Candidate resolution}: a complete candidate proof/disproof with closed
dependencies and formal/certified artifacts, externally reviewed.\\
R9 & \emph{Community validation}: acceptance through the external mathematical
process. \emph{The engine cannot promote itself to this state.}\\
\bottomrule
\end{longtable}
}

Three concurrent tracks feed the ladder. \textbf{Track~A (engineering and
reconstruction)} uses the Bang campaign as the end-to-end integration
laboratory. \textbf{Track~B (blind scientific evaluation)} uses problems supplied
by external mathematicians --- withheld solutions or unpublished subproblems ---
under contamination and information-access controls, which is the only honest way
to measure discovery. \textbf{Track~C (Millennium preparation)} is initially
restricted to constitution, literature atlas, formalized definitions, barriers,
equivalences, certified computational frontiers, and ranked subproblems; it may
advance to direct frontier research only after predetermined gates from Tracks~A
and~B are passed.

\paragraph{Metrics by plane.} The primary trust metric remains the \emph{False
Theorem Promotion Rate} (FTPR): the fraction of \state{PROMOTED} promotions that
are, in fact, wrong. It is never collapsed with discovery metrics into a single
score. Because R2--R9 are open, no promotions exist yet and the FTPR is currently
\emph{undefined}; its estimator is nonetheless specified --- seeded known-false
claims, Track-B blind ground truth, and post-hoc external review supply the
numerator over all \state{PROMOTED} events --- so it becomes measurable exactly
when governance (R2) runs. Around it: (\emph{Discovery}) useful-candidate rate, time-to-counterexample,
novel-result yield, program diversity, reduction depth, theory compression,
barrier discovery, cross-campaign reuse; (\emph{Retrieval/representation})
assumption-compatible retrieval precision, Evidence-Bundle completeness,
definition/version accuracy, dependency-path recall, stale-claim serving rate,
informal--formal alignment accuracy; (\emph{Trust}) FTPR, defect detection and
escape rates, false-rejection rate, node-versus-edge error share, independent
reproduction rate, formalization coverage, re-audit latency, and cost per
correctly promoted claim. A system that solves more problems while raising its
FTPR is, by this architecture's standard, worse.

\section{The Millennium-scale stress test}\label{sec:millennium}

A Millennium Prize Problem is the intended stress test at the top of the ladder,
not a promised deliverable, and the title of this paper makes no such promise.
Its role is to supply adversarial pressure no synthetic benchmark can, at a
target where ``solution'' is defined entirely outside the engine
(Principle~\ref{pr:solution}). Consistent with milestone insurance
(\S\ref{sec:scheduler}), a Millennium campaign is structured so its floor ships
regardless of the frontier: a verified state-of-the-art atlas and formalized
definitions (R6) are a survey \emph{fascicule} that is valuable even if no new
theorem is proved; a frontier result (R7) --- and a proved \otype{BARRIER}
(state \state{NO\_GO}) is valued at parity with a positive result --- is the
aspirational ceiling; and R8--R9 depend on an external process the engine cannot
short-circuit. The ceiling is upside, never the plan.

\section{Threats to validity and limitations}\label{sec:threats}

The engine does not manufacture insight, novelty, tractability, or a Millennium
resolution; a campaign is bounded by the best reduction any Strat\`ege proposes,
and by representation bottlenecks, formalization cost, retrieval error, model
homogenization, benchmark contamination, human-expert bottlenecks, and compute
limits. Against known failure modes, we classify the architecture's response as
\textbf{prevent}, \textbf{detect}, \textbf{record}, \textbf{mitigate}, or
\textbf{open}.

{\footnotesize
\begin{longtable}{@{}p{0.42\textwidth}p{0.16\textwidth}p{0.34\textwidth}@{}}
\caption{Threat model.}\label{tab:threats}\\
\toprule
\textbf{Threat} & \textbf{Response} & \textbf{Mechanism}\\ \midrule
\endfirsthead
\toprule \textbf{Threat} & \textbf{Response} & \textbf{Mechanism}\\ \midrule
\endhead
Hallucinated citations & detect & Biblioth\'ecaire source verification;
retrieval-provenance lens\\
False novelty / renaming & detect & novelty-and-priority lens
(Principle~\ref{pr:novelty})\\
Statement/definition/proof mismatch & prevent & context-hash identity;
definition--proof audit\\
Invalid transformation/reduction edge & detect & edge-level soundness obligation\\
Formalizing a weaker statement & detect & correspondence audit
(Principle~\ref{pr:correspond})\\
Assumption/scope leakage & detect & assumption/scope lens; hard retrieval filters\\
Stale retrieval & prevent & version-aware invalidation of indexes and bundles\\
Corrupted graph edges & record & append-only log; deterministic replay\\
Retrieval poisoning / prompt injection & mitigate & authority membrane; no
promotion from soft memory\\
Checker/searcher common-mode bug & detect & checker-independence lens; mutation
tests\\
Dependency cycles & prevent & acyclicity-of-justification audit\\
Agent collusion / shared-context bias & mitigate & separation of powers;
fresh-context panels\\
Premature convergence to one program & mitigate & hierarchical portfolio;
orthogonal-exploration budget\\
Benchmark contamination & mitigate & Track~B blind problems; contamination
controls\\
Compute-driven publication bias & record & every promotion event logs its
review level and cost\\
Governance capture & open & the top of the ladder is external by construction\\
\bottomrule
\end{longtable}
}

The central limitation is evidential: two modules are implemented prototypes, one
computation is certified, the governance and Congress are so far only manually
exercised, and the discovery, formalization, and scheduling planes are specified,
not built (Table~\ref{tab:maturity}). The immediate work is to climb R2--R4 under
the three-track program and to test the fresh-context hypothesis
(Principle~\ref{pr:fresh}) by varying reviewer independence.

\section{Conclusion}\label{sec:conclusion}

The transition this paper makes is deliberate: a claim-promotion architecture,
preserved intact as the Trust Plane, becomes a governed \emph{discovery}
architecture supported by an implemented representation (\mirador{}) and
retrieval (\pc{}) substrate, and evaluated through readiness gates rather than
assertion. The final position is meant to be ambitious but falsifiable.

\begin{principle}[position]\label{pr:position}
The \engine{} does not claim to manufacture a Millennium solution. It specifies
and evaluates the institutional, representational, computational, and
verification machinery required for autonomous mathematical research to
accumulate frontier progress \emph{without epistemic collapse} --- the state in
which cheap generation silently fills a corpus with load-bearing falsehoods.
\end{principle}

Its reusable asset is horizontal: governed goal graphs, adversarial promotion
boundaries, certificate discipline, and dependency-driven invalidation are a
verified-reasoning infrastructure applicable wherever claims are cheap to
generate and expensive to trust. Each campaign is intended to calibrate the
scheduler's priors and harden the audit checklists, so that --- if the effect
holds --- the engine that finishes a milestone is better calibrated than the one
that started it.

\appendix
\section{Internal language (normative glossary)}\label{app:language}

Bourbaki's working culture survives in its documents, and the engine adopts its
vocabulary because names carry norms. The glossary is normative for platform
documents. \emph{Les \'El\'ements}: the corpus of \state{PROMOTED} nodes with
their proofs and certificates, the only citable ground truth. \emph{Fascicule}:
a shippable closed milestone, the publication quantum. \emph{La Tribu}: the
append-only campaign log, the single source of truth of which graph, corpus, and
indexes are derived views. \emph{Le Congr\`es}: the refutation panel of
\S\ref{sec:congress}. \emph{Le veto}: a structured block on promotion; unanimity
is required at the \state{PROMOTED} boundary. \emph{Tournant dangereux}
(\dbend): a standing flag on nodes or edges with subtle traps --- degenerate
cases, definition drift, normalization choices --- that a panel must explicitly
discharge. \emph{La retraite \`a 50}: mandatory agent retirement --- a context
that has worked a node for $N$ rounds is retired and replaced by a fresh one that
re-derives before extending, institutionalizing fresh eyes. \emph{Le cobaye}: a
deliberately naive agent given a draft cold, to test self-containedness before
the Congress. \emph{Abus de langage}: sanctioned informal shorthand, legal only
when hyperlinked to the rigorous statement it abbreviates. \emph{Plan de
travail}: the Strat\`ege's prioritized frontier snapshot, the only to-do list.

\section{Object types and epistemic states}\label{app:states}

Object type answers ``what is this node?'' and is fixed at creation:
\otype{DEFINITION}, \otype{QUESTION}, \otype{CONJECTURE}, \otype{CLAIM},
\otype{REDUCTION}, \otype{EQUIVALENCE}, \otype{EXAMPLE}, \otype{COUNTEREXAMPLE},
\otype{BARRIER}, \otype{EXPERIMENT}, \otype{CERTIFICATE}, \otype{CHECKER},
\otype{HEURISTIC}, \otype{TRANSLATION}, \otype{PROGRAM}. Epistemic state answers
``how far has it survived?'' and changes only through logged events:
\state{DRAFT} $\to$ \state{OPEN} $\to$ \state{EVIDENCE}/\state{CANDIDATE} $\to$
\state{FORMALIZED} $\to$ \state{AUDITED} $\to$ \state{PROMOTED}, with
\state{REFUTED} and \state{NO\_GO} as terminal-for-a-route outcomes,
\state{SUSPECT} as the only automatically entered state (via invalidation), and
\state{RETIRED} for superseded versions (assigned as part of the promoting event
that mints the new version, not by automatic demotion). The two axes are
independent: a
\otype{CONJECTURE} may be \state{PROMOTED} (as a conjecture, with that status),
and a \otype{CLAIM} labelled ``Theorem'' may sit at \state{DRAFT}
indefinitely (Principle~\ref{pr:orthogonal}).

\begin{thebibliography}{99}
\itemsep2pt

\bibitem{Hilbert1900} D.~Hilbert, \emph{Mathematische Probleme}, Nachr. Ges.
Wiss. G\"ottingen, Math.-Phys. Kl. (1900), 253--297.

\bibitem{Polya45} G.~P\'olya, \emph{How to Solve It}, Princeton University
Press, 1945.

\bibitem{Lakatos76} I.~Lakatos, \emph{Proofs and Refutations}, Cambridge
University Press, 1976.

\bibitem{Bourbaki} N.~Bourbaki, \emph{\'El\'ements de math\'ematique}, Hermann
(later Masson, Springer), Paris, 1939--.

\bibitem{Mashaal06} M.~Mashaal, \emph{Bourbaki: A Secret Society of
Mathematicians}, American Mathematical Society, 2006.

\bibitem{Corry92} L.~Corry, \emph{Nicolas Bourbaki and the concept of
mathematical structure}, Synthese \textbf{92} (1992), 315--348.

\bibitem{Weil92} A.~Weil, \emph{The Apprenticeship of a Mathematician},
Birkh\"auser, 1992.

\bibitem{Beaulieu93} L.~Beaulieu, \emph{A Parisian caf\'e and ten proto-Bourbaki
meetings (1934--1935)}, Math. Intelligencer \textbf{15} (1993), no.~1, 27--35.

\bibitem{NewellSimon57} A.~Newell, J.~C.~Shaw, H.~A.~Simon, \emph{Empirical
explorations of the Logic Theory Machine}, Proc. Western Joint Computer
Conference (1957), 218--230.

\bibitem{Nilsson71} N.~J.~Nilsson, \emph{Problem-Solving Methods in Artificial
Intelligence}, McGraw-Hill, 1971.

\bibitem{Bundy88} A.~Bundy, \emph{The use of explicit plans to guide inductive
proofs}, Proc. CADE-9, LNCS \textbf{310}, Springer (1988), 111--120.

\bibitem{McCune97} W.~McCune, \emph{Solution of the Robbins problem}, J.
Automated Reasoning \textbf{19} (1997), 263--276.

\bibitem{AppelHaken77} K.~Appel, W.~Haken, \emph{Every planar map is four
colorable. Part I: Discharging}, Illinois J. Math. \textbf{21} (1977), 429--490.

\bibitem{Gonthier08} G.~Gonthier, \emph{Formal proof---the four-color theorem},
Notices Amer. Math. Soc. \textbf{55} (2008), 1382--1393.

\bibitem{Hales17} T.~Hales et al., \emph{A formal proof of the Kepler
conjecture}, Forum of Mathematics, Pi \textbf{5} (2017), e2.

\bibitem{Heule16} M.~J.~H.~Heule, O.~Kullmann, V.~W.~Marek, \emph{Solving and
verifying the Boolean Pythagorean triples problem via cube-and-conquer}, Proc.
SAT 2016, LNCS \textbf{9710}, Springer (2016), 228--245.

\bibitem{mathlib20} The mathlib Community, \emph{The Lean mathematical library},
Proc. 9th ACM SIGPLAN Conf. Certified Programs and Proofs (CPP 2020), 367--381.

\bibitem{deBruijn70} N.~G.~de Bruijn, \emph{The mathematical language AUTOMATH},
Symposium on Automatic Demonstration, LNM \textbf{125}, Springer (1970), 29--61.

\bibitem{GowersNielsen09} T.~Gowers, M.~Nielsen, \emph{Massively collaborative
mathematics}, Nature \textbf{461} (2009), 879--881.

\bibitem{Polymath12} D.~H.~J.~Polymath, \emph{A new proof of the density
Hales--Jewett theorem}, Ann. of Math. \textbf{175} (2012), 1283--1327.

\bibitem{Yao23} S.~Yao et al., \emph{Tree of Thoughts: deliberate problem
solving with large language models}, NeurIPS 36 (2023).

\bibitem{Trinh24} T.~H.~Trinh, Y.~Wu, Q.~V.~Le, H.~He, T.~Luong, \emph{Solving
olympiad geometry without human demonstrations}, Nature \textbf{625} (2024),
476--482.

\bibitem{RomeraParedes24} B.~Romera-Paredes et al., \emph{Mathematical
discoveries from program search with large language models}, Nature \textbf{625}
(2024), 468--475.

\bibitem{Davies21} A.~Davies et al., \emph{Advancing mathematics by guiding
human intuition with AI}, Nature \textbf{600} (2021), 70--74.

\bibitem{Aletheia26} T.~Feng, T.~H.~Trinh, G.~Bingham, et al.\ (Google
DeepMind), \emph{Towards Autonomous Mathematics Research} (the Aletheia
system), arXiv:2602.10177 (2026).

\bibitem{QED26} C.~An, Q.~Ye, M.~Pan, J.~Zhang, \emph{QED: an open-source
multi-agent system for generating mathematical proofs on open problems},
arXiv:2604.24021 (2026).

\bibitem{RMA26} Z.~Zhao, B.~Yuan, J.~Choi, Y.~Chen, \emph{RMA: an agentic system
for research-level mathematical problems}, arXiv:2605.22875 (2026).

\bibitem{ProverAgent25} Baba, Liu, Kurita, Sannai, \emph{Prover Agent: an
agent-based framework for formal mathematical proofs}, arXiv:2506.19923 (2025).

\bibitem{LEAP26} Kung, Song, Hwang, et~al., \emph{LEAP: supercharging LLMs for
formal mathematics with agentic frameworks}, arXiv:2606.03303 (2026).

\bibitem{LeanSearch26} Gao, Sun, Jiang, et~al., \emph{LeanSearch v2: global
premise retrieval for Lean~4 theorem proving}, arXiv:2605.13137 (2026).

\bibitem{AlphaEvolve25} Novikov, Balog, et~al.\ (Google DeepMind),
\emph{AlphaEvolve: a coding agent for scientific and algorithmic discovery},
Google DeepMind technical report (2025),
\texttt{deepmind.google/discover/blog/alphaevolve}.

\bibitem{Doyle79} J.~Doyle, \emph{A truth maintenance system}, Artificial
Intelligence \textbf{12} (1979), no.~3, 231--272.

\bibitem{deKleer86} J.~de~Kleer, \emph{An assumption-based TMS}, Artificial
Intelligence \textbf{28} (1986), 127--162.

\bibitem{Celik17} A.~Celik, K.~Palmskog, M.~Gligoric, \emph{iCoq: regression
proof selection for large-scale verification projects}, Proc. ASE 2017.

\bibitem{Massot} P.~Massot, \emph{leanblueprint: a plasTeX plugin for writing
formalization blueprints for Lean~4 projects}, software.

\bibitem{Tao23} T.~Tao, \emph{Formalizing the proof of PFR in Lean4 using
Blueprint: a short tour}, blog post, November 2023.

\bibitem{vanDoorn20} F.~van~Doorn, G.~Ebner, R.~Y.~Lewis, \emph{Maintaining a
library of formal mathematics}, Proc. CICM 2020, LNCS \textbf{12236}, Springer
(2020); arXiv:2004.03673.

\bibitem{TheoremGraph26} S.~Kurgan, E.~Wang, E.~Leonen, et al., \emph{TheoremGraph:
bridging formal and informal mathematics}, arXiv:2606.25363 (2026).

\bibitem{MaRDI23} M.~Schubotz, E.~Ferrer, J.~Stegm\"uller, et al., \emph{Bravo
MaRDI: a Wikibase powered knowledge graph on mathematics}, arXiv:2309.11484
(2023).

\bibitem{PROVDM13} L.~Moreau, P.~Missier (eds.), \emph{PROV-DM: the PROV data
model}, W3C Recommendation, 30 April 2013.

\bibitem{Kuhn14} T.~Kuhn, M.~Dumontier, \emph{Trusty URIs: verifiable,
immutable, and permanent digital artifacts for linked data}, Proc. ESWC 2014,
LNCS \textbf{8465}, Springer (2014), 395--410.

\bibitem{Kuhn16} T.~Kuhn et al., \emph{Decentralized provenance-aware publishing
with nanopublications}, PeerJ Computer Science \textbf{2} (2016), e78.

\bibitem{RFC6962} B.~Laurie, A.~Langley, E.~Kasper, \emph{Certificate
transparency}, RFC 6962, IETF, 2013.

\bibitem{Cox19} R.~Cox, \emph{Transparent logs for skeptical clients},
\texttt{research.swtch.com/tlog}, 2019.

\bibitem{Fowler05} M.~Fowler, \emph{Event sourcing}, \texttt{martinfowler.com},
2005.

\bibitem{Zep25} P.~Rasmussen, P.~Paliychuk, T.~Beauvais, J.~Ryan, D.~Chalef,
\emph{Zep: a temporal knowledge graph architecture for agent memory},
arXiv:2501.13956 (2025).

\bibitem{Yadav26} N.~Yadav, \emph{Temporal validity in retrieval memory},
arXiv:2606.26511 (2026).

\bibitem{Auer02} P.~Auer, N.~Cesa-Bianchi, P.~Fischer, \emph{Finite-time
analysis of the multiarmed bandit problem}, Machine Learning \textbf{47} (2002),
235--256.

\bibitem{Browne12} C.~Browne et al., \emph{A survey of Monte Carlo tree search
methods}, IEEE Trans. Comput. Intell. AI in Games \textbf{4} (2012), 1--43.

\bibitem{Bang51} T.~Bang, \emph{A solution of the ``plank problem''}, Proc.
Amer. Math. Soc. \textbf{2} (1951), 990--993.

\bibitem{Ball91} K.~Ball, \emph{The plank problem for symmetric bodies}, Invent.
Math. \textbf{104} (1991), 535--543.

\bibitem{Gardner88} R.~J.~Gardner, \emph{Relative width measures and the plank
problem}, Pacific J. Math. \textbf{135} (1988), 299--312.

\bibitem{AKP19} A.~Akopyan, R.~Karasev, F.~Petrov, \emph{Bang's problem and
symplectic invariants}, J. Symplectic Geom. \textbf{17} (2019), no.~6,
1579--1611.

\bibitem{GKP22} A.~Glazyrin, R.~Karasev, A.~Polyanskii, \emph{Covering by planks
and avoiding zeros of polynomials}, Int. Math. Res. Not. IMRN (2023);
arXiv:2112.05382.

\bibitem{Ambrus22} G.~Ambrus, \emph{A short proof of the plank theorem for
symmetric bodies}, arXiv:2203.11260 (2022); withdrawn by the author.

\bibitem{Verreault26} W.~Verreault, \emph{Plank theorems and their applications:
a survey}, Bull. London Math. Soc. (2026); arXiv:2203.05540.

\bibitem{BY26} A.~Bakaev, A.~Yehudayoff, \emph{A lower bound for the total
relative width of planks covering a convex body}, arXiv:2602.20290 (2026).

\bibitem{MOM26} A.~D.~Mart\'inez, O.~Ortega-Moreno, \emph{A solution to the
strong polarization problem}, arXiv:2605.28744 (2026).

\bibitem{Lucius26} M.~Lucius, \emph{Transport and tiling bounds for the affine
plank problem on the triangle}, manuscript (2026), 35 pp.; repository
\texttt{maximilianolucius/bang-plank-affine-triangle}.

\bibitem{lucius2026mirador} M.~Lucius, \emph{MIRADOR: a typed, versioned
intermediate representation for auditable mathematical knowledge and discovery},
software and preprint (2026), Zenodo \texttt{10.5281/zenodo.21324524};
repository \texttt{maximilianolucius/mirador}.

\bibitem{lucius2026proofcontext} M.~Lucius, \emph{ProofContext: dependency-aware
retrieval of verifiable context for mathematical reasoning}, software and
preprint (2026), Zenodo \texttt{10.5281/zenodo.21325288}; repository
\texttt{maximilianolucius/proofcontext}.

\end{thebibliography}

\end{document}
